\begin{lemma}
\label{lemma-indecomposable-injective}
Let $R$ be a ring. Let $E$ be an indecomposable injective $R$-module.
Then
\begin{enumerate}
\item $E$ is the injective hull of any nonzero submodule of $E$,
\item the intersection of any two nonzero submodules of $E$ is nonzero,
\item $\text{End}_R(E)$ is a noncommutative local ring with maximal
ideal those $\varphi : E \to E$ whose kernel is nonzero, and
\item the set of zerodivisors on $E$ is a prime ideal $\mathfrak p$ of $R$
and $E$ is an injective $R_\mathfrak p$-module.
\end{enumerate}
\end{lemma}

\begin{proof}
Part (1) follows from Lemma \ref{lemma-injective-hull-unique}.
Part (2) follows from part (1) and the definition of injective hulls.

\medskip\noindent
Proof of (3). Set $A = \text{End}_R(E)$ and
$I = \{\varphi \in A \mid \Ker(\varphi) \not = 0\}$.
The statement means that $I$ is a two sided ideal and
that any $\varphi \in A$, $\varphi \not \in I$ is invertible.
Suppose $\varphi$ and $\psi$ are not injective.
Then $\Ker(\varphi) \cap \Ker(\psi)$ is nonzero
by (2). Hence $\varphi + \psi \in I$. It follows that $I$
is a two sided ideal. If $\varphi \in A$, $\varphi \not \in I$,
then $E \cong \varphi(E) \subset E$ is an injective submodule,
hence $E = \varphi(E)$ because $E$ is indecomposable.

\medskip\noindent
Proof of (4). Consider the ring map $R \to A$ and let $\mathfrak p \subset R$
be the inverse image of the maximal ideal $I$. Then it is clear
that $\mathfrak p$ is a prime ideal and that $R \to A$ extends to
$R_\mathfrak p \to A$. Thus $E$ is an $R_\mathfrak p$-module.
It follows from Lemma \ref{lemma-injective-epimorphism} that $E$ is injective
as an $R_\mathfrak p$-module.
\end{proof}
